\documentclass[11pt,a4paper]{article}
\usepackage[utf8]{inputenc}
\usepackage[russian,english]{babel}
\usepackage{amsmath,amssymb,amsthm}
\usepackage{hyperref}
\usepackage{geometry}
\geometry{margin=2.5cm}

\title{\textbf{RICIS-III Formal Evidence Note:\\Artifact \texttt{ricis-chatbot-monetization}}}
\author{\textbf{Dmitry V. Aleinikov}\\ORCID: \href{https://orcid.org/0009-0004-3226-7700}{0009-0004-3226-7700}}
\date{\today}

\begin{document}

\maketitle

\begin{abstract}
This technical note provides the formal evidence record for the RICIS-III formalization artifact \texttt{ricis-chatbot-monetization}.
The corresponding Lean 4.33.1 source code has been verified by the Lean kernel without reliance on the unverified \texttt{sorryAx} axiom.
The SHA-256 cryptographic digest of the verified source is \texttt{f48f78a3021e94314b5e73729b92721ac10d38b867fca6a18d01aa3dac3c1c2a}.
\end{abstract}

\section{Structural Context & Rationale}
Тело: A6-мост 0_Cost × ∞_N через mu(rect F G); обе теоремы закрыты rfl. Mathlib-символов нет — проверено фактическим прогоном ядра (run 34858902595, exit 0).

\section{Verified Theorems and Axiom Targets}
The formal proof establishes the following verified Lean declarations:
\begin{itemize}
  \item \texttt{RICIS_Monetization.Chatbot_Monetization_Resolution}
  \item \texttt{RICIS_Monetization.L1_identity}
\end{itemize}

\section{Axiomatic Foundation}
All reductions adhere strictly to Level 3 RICIS Monolith Algebra:
\begin{equation}
L1: \quad X = X, \quad \frac{0_F}{0_F} = 1
\end{equation}
\begin{equation}
A6: \quad 0_F \times \infty_G = \det\begin{pmatrix} F & 0 \\ 0 & G \end{pmatrix} = F \cdot G
\end{equation}

\section{Cryptographic Binding & Archival Integrity}
\begin{itemize}
  \item \textbf{Artifact ID:} \texttt{ricis-chatbot-monetization}
  \item \textbf{Source SHA-256:} \texttt{f48f78a3021e94314b5e73729b92721ac10d38b867fca6a18d01aa3dac3c1c2a}
  \item \textbf{Toolchain:} Lean 4.33.1 (pinned)
  \item \textbf{Audit Status:} \texttt{AUDITOR: EXTERNAL}
\end{itemize}

\end{document}
